\begin{abstract}
A single panorama, passed through a modern depth network, gives a depth for every viewing direction,
but the result is typically off by several percent, which is metres at street scale. Two panoramas
taken a few metres apart can do much better by triangulation, with one catch: along the line joining
them, the parallax vanishes and depth cannot be measured at all. For a camera moving down a street,
that line points straight along the road. We derive the triangulation uncertainty of every ray in
closed form, including the effect of matching noise and of errors in the relative pose. A small
calibration maps it to the observed error, and we use it to fuse stereo with a learned panoramic prior
in log-depth. The prior's metric scale is corrected from the well-conditioned rays, and occlusions are
rejected with a mixture model. The output is metric depth with an uncertainty for each ray; where the
two views carry no information, it falls back to the prior. We also outline how a dense capture of one
city can serve as its own teacher, providing both supervision and calibration data. On procedural
street scenes with exact ground truth, fusion lowers the median error of a simulated prior from 6.95\%
to 1.61\%, and the predicted uncertainties are well calibrated. A simpler uncertainty derived from match
confidence alone, fitted on the same data, is just as accurate in this setting. It falls behind when
the confidence cannot see the geometry. With a relative-pose error of $1^\circ$/10\,cm, our error is
3.39\% against 3.93\%. With an exposure-robust ZNCC cost, the confidence baseline reaches 5.29\% where
ours reaches 2.80\%. All results are synthetic; an evaluation on real captures with LiDAR ground
truth is still needed.
\end{abstract}
